% Relazione di progetto - Laboratorio di Programmazione di Sistemi Mobili
% Compilare con: pdflatex relazione.tex
\documentclass[11pt,a4paper]{article}

\usepackage[utf8]{inputenc}
\usepackage[T1]{fontenc}
\usepackage[italian]{babel}
\usepackage{lmodern}
\usepackage{microtype}
\usepackage[top=2.8cm,bottom=2.8cm,left=3.2cm,right=3.2cm]{geometry}
\usepackage{parskip}
\usepackage{booktabs}
\usepackage{listings}
\usepackage{xcolor}
\usepackage{hyperref}
\hypersetup{colorlinks=true, linkcolor=black, urlcolor=blue!55!black}

% un po' d'aria fra le righe e attorno ai titoli
\linespread{1.12}
\setlength{\parskip}{0.7\baselineskip}
\setlength{\emergencystretch}{1em}

\definecolor{codebg}{gray}{0.96}
\lstset{
  basicstyle=\ttfamily\small,
  backgroundcolor=\color{codebg},
  breaklines=true,
  columns=fullflexible,
  keepspaces=true,
  showstringspaces=false,
  xleftmargin=1em,
  framexleftmargin=1em,
  framexrightmargin=1em,
  framextopmargin=6pt,
  framexbottommargin=6pt,
  frame=tb,
  framerule=0pt,
  aboveskip=1.2em,
  belowskip=1.2em,
}

\begin{document}

\thispagestyle{empty}

\begin{center}
  \vspace*{1.5cm}

  {\LARGE\bfseries TaskManager}

  \vspace{0.6em}

  {\large Relazione di progetto}

  \vspace{2.2em}

  % TODO: mettere i propri dati
  Nome Cognome \quad--\quad matricola 000000

  \texttt{nome.cognome@studio.unibo.it}

  \vspace{2em}

  Laboratorio di Programmazione di Sistemi Mobili\\
  a.a.\ 2025/26 \quad--\quad progetto individuale, Android e Kotlin

  \vspace{1em}

  \today
\end{center}

\vspace{2cm}


\section*{Da dove nasce}
\addcontentsline{toc}{section}{Da dove nasce}

Per un altro corso mi ero fatto una web app per tenere in ordine i task.
Titolo, priorità, stato, due tag, una scadenza. Niente di che, però la usavo.

Il problema è che girava nel browser. E i task uno se li ricorda in fila alla
posta, non seduto davanti al portatile. Da lì l'idea: rifarla su Android.

Il backend l'ho lasciato dov'era, su Supabase, con le stesse tabelle. Lo dico
subito perché sembra una scorciatoia, e un po' lo è. Ma avere il database già in
piedi mi ha lasciato tutto il tempo per la parte Android, che poi è quella
dell'esame. Il conto l'ho pagato da un'altra parte: certe scelte di Supabase me
le sono trovate addosso, e un paio di serate le ho buttate lì dentro. Ne parlo
più avanti, perché è la parte in cui ho imparato qualcosa.

Cosa fa l'app, in breve. Ti registri, fai login. Vedi i tuoi task come card: il
colore è la priorità, lo stato è un'etichetta che avanza se ci tocchi sopra,
senza aprire niente. Crei, modifichi, elimini. Tag, scadenza con data e ora.
Sopra c'è una barra con ricerca, filtri e ordinamento. Chiudi l'app, la riapri,
sei ancora dentro. C'è l'italiano e c'è l'inglese.

Una cosa da sapere prima di leggere il codice: lato server non ho scritto
niente. Le quattro tabelle stanno su Supabase ed escono via REST. Chi vede cosa
lo decide il database, con le policy di Row Level Security. Ecco perché nelle
GET non filtro mai per \texttt{user\_id}: non serve, ci pensa lui guardando il
token con cui arrivo.


\section*{Come è fatto dentro}
\addcontentsline{toc}{section}{Come è fatto dentro}

Quattro moduli Gradle. E qui mi fermo un attimo, perché la divisione in moduli
sembra sempre una cosa messa lì per la relazione, e invece a me serviva davvero.

Le dipendenze fra i moduli vanno in un verso solo. Risultato: se provo a
scrivere una chiamata di rete dentro una schermata, non compila. Non è una
regola che devo ricordarmi io, è una regola che me la ricorda Gradle.

\begin{table}[h]
\centering
\begin{tabular}{@{}llp{6.2cm}@{}}
\toprule
\textbf{Modulo} & \textbf{Dipende da} & \textbf{Cosa c'è dentro} \\
\midrule
\texttt{:domain} & niente & modelli, interfacce dei repository, use case \\
\addlinespace[0.3em]
\texttt{:data} & \texttt{:domain} & Retrofit, Moshi, DataStore, i repository veri \\
\addlinespace[0.3em]
\texttt{:uicompose} & \texttt{:domain} & schermate Compose e ViewModel \\
\addlinespace[0.3em]
\texttt{:app} & tutti & \texttt{Application}, manifest, cablaggio \\
\bottomrule
\end{tabular}
\end{table}

\texttt{:domain} è il cuore, ed è povero apposta. Non dipende da nessuno.
L'unica libreria che si porta dietro sono le coroutine, e solo perché i
repository espongono dei \texttt{Flow}. Dentro ci sono i modelli, le interfacce
dei due repository (solo le interfacce: il contratto) e gli use case, uno per
file.

\texttt{:data} è dove finisce tutto quello che tocca l'esterno. Le due
interfacce Retrofit, i DTO di Moshi con i mapper, il \texttt{SessionManager} che
tiene la sessione su DataStore, e i repository veri. Le chiavi di Supabase le
tengo in \texttt{local.properties}, che non committo; arrivano al codice
passando per il \texttt{BuildConfig}.

\texttt{:uicompose} sono le tre schermate: login, dashboard, form. Ognuna col
suo ViewModel e la sua factory. In \texttt{common} ci sono i pezzi che riuso,
badge, banner d'errore, loader, la formattazione delle date. Questo modulo non
sa nemmeno che esista una rete.

\texttt{:app} è quattro file in croce: manifest, tema, e la
\texttt{CustomApplication}. Che però fa la cosa che regge tutto il resto.

\begin{lstlisting}
class CustomApplication : Application() {
    override fun onCreate() {
        super.onCreate()
        UseCasesProvider.setup(
            repositoryProvider = RepositoryProviderImpl(this)
        )
    }
}
\end{lstlisting}

Dependency injection a mano, senza Hilt, come a lezione. Mi piace che il punto
in cui si incastra tutto siano cinque righe leggibili. Mi piace meno che per
aggiungere una dipendenza debba aprire tre file. Per un progetto di questa
taglia, però, va benissimo così.


\section*{I requisiti d'esame, dove stanno}
\addcontentsline{toc}{section}{I requisiti d'esame, dove stanno}

\begin{table}[h]
\centering
\begin{tabular}{@{}p{4.2cm}p{9.4cm}@{}}
\toprule
\textbf{Requisito} & \textbf{Dove sta} \\
\midrule
Separazione in package e moduli &
Quattro moduli Gradle, dipendenze in un verso solo, package divisi per
responsabilità dentro a ciascuno. \\
\addlinespace[0.5em]
Componenti di lifecycle &
Un \texttt{ViewModel} per schermata, con la sua factory. Lo stato esce come
\texttt{StateFlow} e la UI lo raccoglie con
\texttt{collectAsStateWithLifecycle}. Quel poco che vive nella schermata sta in
\texttt{rememberSaveable}, così regge la rotazione. \\
\addlinespace[0.5em]
Coroutine, main thread separato dal background &
Repository in \texttt{withContext(Dispatchers.IO)}, ViewModel in
\texttt{viewModelScope}. La UI resta sul main thread e riceve solo flussi già
pronti. Le quattro GET del refresh partono insieme con \texttt{async}, visto che
non dipendono l'una dall'altra. \\
\addlinespace[0.5em]
Almeno due chiamate ad API remote &
Dieci endpoint. \texttt{signup} e \texttt{token} (login e rinnovo) per
l'autenticazione; \texttt{GET}, \texttt{POST}, \texttt{PATCH} e \texttt{DELETE}
su \texttt{tasks}; una \texttt{GET} a testa su \texttt{tags},
\texttt{priorities} e \texttt{statuses}. \\
\addlinespace[0.5em]
Opzionale 1: salvataggio in uno storage locale &
Il \texttt{SessionManager} scrive access token, refresh token, id utente e
scadenza su DataStore Preferences, e li rilegge all'avvio: l'app riparte già
autenticata. Al logout svuota tutto. \\
\bottomrule
\end{tabular}
\end{table}


\section*{Le cose di cui sono contento}
\addcontentsline{toc}{section}{Le cose di cui sono contento}

\subsection*{Le dipendenze puntano tutte verso il centro}

Il vincolo che ho tenuto più stretto è che \texttt{:uicompose} non vede
\texttt{:data}. Non è buona volontà: un ViewModel non riesce proprio a scrivere
il nome di \texttt{TaskRepositoryImpl}, perché quella classe non ce l'ha nel
classpath. Tutto passa dagli use case, che parlano solo di interfacce dichiarate
in \texttt{:domain}.

Se domani volessi buttare Supabase e mettere altro, riscriverei \texttt{:data} e
non toccherei una riga di interfaccia. Stessa cosa per la cache locale che non
ho fatto in tempo ad aggiungere: entrerebbe tutta lì dentro, dietro alle stesse
interfacce, e le schermate non se ne accorgerebbero.

\subsection*{La serata persa dietro a un 200 OK}

Questa la racconto tutta, perché è la sera in cui ho capito più cose.

Riavvio il telefono, apro l'app: zero task. Nessun crash. Nessun errore. Nessuna
riga rossa nel log. Solo una lista vuota, come se non avessi mai creato niente.
E il token stava lì, lo vedevo dentro DataStore.

Ci ho messo un po' a trovarlo. Il fatto è che leggere da DataStore è asincrono.
Le chiamate partivano prima che il token tornasse dal disco, quindi partivano
senza l'header \texttt{Authorization}. Fin qui ci sta.

La parte cattiva è la risposta. Supabase non ti dice che non sei autenticato. La
Row Level Security fa il suo lavoro, non trova righe che io possa vedere, e
risponde \texttt{200 OK} con una lista vuota. Un errore di autenticazione
vestito da risposta giusta.

Da lì la soluzione. Il \texttt{SessionManager} alza un flag quando ha finito di
leggere il disco, e ogni chiamata di rete lo aspetta prima di partire. Stessa
attesa dentro \texttt{AuthRepositoryImpl}: \texttt{isLoggedIn} sta zitto finché
il ripristino non è concluso. Se no, per un istante la UI vede un falso ``non
loggato'' e sbatte la schermata di login in faccia a uno che era già dentro.

E già che ci sono: è lo stesso motivo per cui ho due client HTTP invece di uno.
Gli endpoint di login vogliono solo la chiave anonima del progetto, quelli sulle
tabelle vogliono anche il token dell'utente come \texttt{Bearer}. Sembra roba
duplicata, finché non mandi la chiave anonima sulle tabelle e ti ritrovi,
indovina un po', un altro \texttt{200 OK} con la lista vuota.

\subsection*{Gli errori sono tipi, non frasi}

Il modulo \texttt{:data} non sa in che lingua parla l'utente. Quindi non è lui
che deve scrivere i messaggi. Espone tre casi e basta:

\begin{lstlisting}
sealed class ErrorState {
    data object Network : ErrorState()
    data object Unauthorized : ErrorState()
    data object Generic : ErrorState()
}
\end{lstlisting}

Il testo lo mette la UI, pescando da \texttt{values/strings.xml} o da
\texttt{values-it/strings.xml}. Se domani serve il francese aggiungo un file di
stringhe e non tocco nient'altro.

\subsection*{Filtrare non deve chiamare la rete}

Filtri, ricerca e ordinamento sono una funzione pura, in \texttt{:domain}. Le
dai la lista e i filtri, ti torna la lista da mostrare. Il
\texttt{HomeViewModel} la richiama con \texttt{combine} ogni volta che cambia
qualcosa. Quindi scrivere nella casella di ricerca non fa partire nessuna
chiamata.

E siccome quella funzione non sa niente di Android, la provo con una JVM
normale, senza emulatore acceso. Che poi è il criterio con cui ho scelto dove
mettere i test: dove il codice si può rompere in silenzio. Il parsing delle
risposte di login, compreso il caso in cui la registrazione chiede la conferma
via mail e non torna nessun token. Il mapping delle righe di PostgREST, con i
campi che mancano e \texttt{tag\_ids} nullo. E i filtri. Per riuscire a testare
i mapper li ho resi \texttt{internal} invece che privati dentro al repository:
ci ho rimesso un po' di incapsulamento, ci ho guadagnato i test dove servivano.

\subsection*{Meglio zoppicare che morire}

Ci sono tre punti in cui l'app poteva crashare per colpa di roba che non dipende
da me, e in tutti e tre ho preferito farla zoppicare.

Una riga che arriva senza \texttt{id} la butto via con \texttt{mapNotNull},
invece di far saltare il parsing di tutta la lista. Se DataStore è illeggibile
risulti non loggato: scomodo, ma vivo. E se \texttt{local.properties} manca del
tutto, l'URL di base diventa un host palesemente finto, così l'app parte lo
stesso e le chiamate falliscono con un normale errore di rete, invece di
piantarsi all'avvio.


\section*{Cosa manca}
\addcontentsline{toc}{section}{Cosa manca}

Room, prima di tutto. Oggi senza rete l'app non mostra niente, ed è il limite
che si sente di più usandola. Il posto però c'è già: la dipendenza è dichiarata,
\texttt{TaskRepositoryImpl} ha già il \texttt{Context}, e
\texttt{startFetchTasks()} l'ho tenuto separato da \texttt{refreshTasks()}
proprio pensando a quel giorno. Mancano entity, DAO e database. Quando ci sarà,
va svuotata al logout, se no due account sullo stesso telefono si vedono i dati
a vicenda.

Poi la foto sul task. Il campo \texttt{photoUrl} nel modello c'è, il permesso
\texttt{CAMERA} nel manifest c'è, il resto no: permesso a runtime, scatto,
caricamento su Supabase Storage, e la colonna \texttt{photo\_url} sulla tabella,
che oggi non esiste. È per quello che il campo non lo mando in scrittura: la
richiesta fallirebbe.

Terzo, un worker che chiami \texttt{RefreshTasksUseCase} ogni tanto, o quando
torna la linea. Lo use case l'ho tenuto isolato apposta.

Poi ci sono le cose meno urgenti. La navigazione, per dire: oggi è una variabile
dentro \texttt{MainActivity}, dove \texttt{null} vuol dire dashboard, stringa
vuota vuol dire ``sto creando'' e un id vuol dire ``sto modificando quello''.
Funziona ed è pochissimo codice. Alla quarta schermata non si capirebbe più
niente e passerei a Navigation Compose. Nella lista lunga ci sono anche il
riordino trascinando le card (la colonna \texttt{sort\_order} c'è già), la
paginazione quando i task diventano tanti, e qualche test sulle schermate.


\section*{Per compilarlo}
\addcontentsline{toc}{section}{Per compilarlo}

Serve un \texttt{local.properties} nella cartella radice, con due righe:

\begin{lstlisting}
SUPABASE_URL=https://<progetto>.supabase.co
SUPABASE_ANON_KEY=<chiave anonima>
\end{lstlisting}

Il file non è versionato, per ovvi motivi. Il \texttt{minSdk} è 31, cioè Android
12, così ho \texttt{java.time} senza doverlo ritradurre con il desugaring;
\texttt{compileSdk} e \texttt{targetSdk} sono 36. I test si lanciano con
\texttt{./gradlew test} e non vogliono nessun emulatore acceso.

\end{document}
